\SourceLeaf{15}
\[
-\frac{\mathrm{d}s'.\Phi r'.\cos MQ'}{r'r'}-\frac{\mathrm{d}s''.\Phi r''.\cos MQ''}{r''r''}-\frac{\mathrm{d}s'''.\Phi r'''.\cos MQ'''}{r'''r'''}-\text{ etc.}
\]
\par
und im zweiten Fall durch dieselbe Formel mit dem zusätzlichen Glied
\[
-\mathrm{d}\Sigma.\Phi0
\]
\par
Multiplizieren wir diesen Ausdruck mit $\cos MX$, so erhalten wir die Kraft, mit der die innerhalb des Kegels liegenden Teile des Körpers den Punkt in der zur $x$-Koordinatenachse parallelen und ihr entgegengesetzten Richtung anziehen. Die Kraft, mit der der ganze Körper in derselben Richtung wirkt, wird daher durch das über die ganze Oberfläche des Körpers erstreckte Integral $-\int\frac{\mathrm{d}s.\Phi r.\cos MQ.\cos MX}{rr}$ ausgedrückt, sofern der angezogene Punkt außerhalb des Körpers liegt. Doch ist auch das über die ganze Kugeloberfläche erstreckte Integral $-\Phi0.\int\mathrm{d}\Sigma.\cos MX$ hinzuzufügen, wenn $M$ innerhalb des Körpers liegt. Man sieht ferner leicht, dass, wenn $M$ auf der Oberfläche des Körpers liegt, zwar dasselbe Integral $-\Phi0.\int\mathrm{d}\Sigma.\cos MX$ hinzuzufügen ist, aber nur über die halbe Kugeloberfläche erstreckt. Es wird über die Halbkugelfläche erstreckt, die durch die Tangentialebene an die Oberfläche des Körpers in $M$ begrenzt wird und auf derselben Seite der Ebene wie der feste Teil des Körpers beim Punkt $M$ liegt. Um den Wert dieses Integrals zu bestimmen, denken wir uns den von dieser Halbkugelfläche und der Ebene eingeschlossenen Körper. Es bezeichne $\theta$ allgemein den Winkel zwischen einer nach außen gerichteten Normalen zur Oberfläche dieses Körpers und einer Geraden parallel zur $x$-Koordinatenachse. Nach dem ersten Lehrsatz verschwindet das über die \textit{ganze} Oberfläche des Körpers erstreckte Integral $\int\mathrm{d}s.\cos\theta$. Setzt man daher das nur über den ebenen Teil der Oberfläche erstreckte Integral $=J$, so muss das über den gekrümmten Teil der Oberfläche erstreckte Integral $=-J$ sein. Auf dem gekrümmten Teil stimmt aber $\mathrm{d}s$ mit unserem $\mathrm{d}\Sigma$ überein, und $\theta$ ist $=180^0-MX$. Für das über die Halbkugelfläche erstreckte Integral $-\int\mathrm{d}\Sigma.\cos MX$ erhält man daher $=-J$. Auf dem ebenen Teil der Oberfläche ist $\theta$ dagegen offenbar konstant und gleich dem Wert von $QX$ im Punkt $M$. Also ist $J$ gleich dem Produkt aus dem Kosinus dieses Winkels und dem Flächeninhalt des ebenen Stücks, der $=\pi$ ist. Daraus ergibt sich für das über die oben bestimmte Halbkugelfläche erstreckte Integral $-\Phi0.\int\mathrm{d}\Sigma.\cos MX$ der Wert $=-\pi\Phi0.\cos QX$, wobei für $QX$ der Wert im Punkt $M$ zu nehmen ist. Ganz ebenso findet man für das über die andere Halbkugelfläche erstreckte Integral $-\Phi0.\int\mathrm{d}\Sigma.\cos MX$ den Wert $=+\pi\Phi0.\cos QX$, sodass das Integral über die ganze Kugel $=0$ ist. Aus all dem erhalten wir
